\section{Uniform selection of the reader comparison on a plateau}
\label{sec:actual-reader-dawn-selection}

We now put the local estimates on one selected history. The choices
must be uniform in the finite period, the reader patch, the forest
and the tag. A further distinction is essential: the temperature used
to record a Gaussian sector need not be the physical coupling at
which the field is sampled. We keep these parameters separate.

\subsection{One allowance for the whole coupling band}

Fix a blocking patch with side $L$, its predictor kernel and adjacent
path data. Fix positive strip width $s$ and lower metric bound $a$,
an upper price $z\ge0$, a reader-room multiplier $w\in\mathbb N$,
a real error exponent $f$, and a real coupling-drift bound $\gamma$.
The preselected block side must satisfy the predictor smallness
condition
\[
 2\dim(\mathfrak g)\,T_\Psi(z,L)\le 1/4096,
\]
where $T_\Psi(z,L)$ is the absolute-row ticket for the truncated
coarse-to-fine curvature transport. This is a condition on the
previous choice of $L$, not something achieved by increasing the
entrance coupling afterward.

For $t>1$ define
\[
 \ell(t)=\log\log t,\qquad
 \varepsilon_f(t)=\exp\{-f\ell(t)^2\}.
\]
For every fixed $v>0$ and real $c$, eventually in $t$,
\begin{equation}\label{eq:reader-band-power}
 x^{-v}\le\varepsilon_c(t)\qquad(t/2\le x\le t).
\end{equation}
In fact $x^{-v}\le 2^vt^{-v}$, and
$c(\log\log t)^2=o(\log t)$. Taking logarithms proves the
inequality, uniformly over the displayed interval. This elementary
comparison lets us use one error allowance at the top of a plateau.

Here are the local inputs to which we apply it. The local Gaussian
comparison gives one entry bound $U\ge1$, followed by an eventual
threshold for its mean and covariance estimates. The nonlinear
guarded comparison supplies its field estimates for all
$x/2\le\beta\le3x$, and the guard bound gives lost mass at most
$x^{-1}$. The common covariance-cover and transport estimates
supply both short-row errors at most $x^{-1/8}$. The adjacent
coupling comparison holds uniformly for
\begin{equation}\label{eq:reader-coupling-band}
 x/2\le\beta\le3x,\qquad x/2\le\beta',\qquad
 |\beta'-\beta|\le\gamma.
\end{equation}
There is no independent upper restriction on $\beta'$. Finally,
the local chart thresholds give the room and logarithm inequalities
used in the next subsection. These are respectively the conclusions
of Sections~\ref{sec:actual-reader-local-gaussian-prices},
\ref{sec:actual-reader-guarded-curvature} and
\ref{sec:actual-reader-short-covariance}, with their stated
same-data and prior-purchase hypotheses. Their production is a
prerequisite here; the following choice argument does not supply an
estimate absent from one of them.

Take the maximum of their finitely many thresholds, after $L$ and
the other displayed fixed parameters have been chosen. For all
sufficiently large $t$, every $x\in[t/2,t]$ is above that maximum.
Set
\[
 p_x=x^{-1/8}.
\]
This value is chosen before $\beta,\beta'$ and pays the short-row
errors for every coupling pair in \eqref{eq:reader-coupling-band}.
Increase the common threshold once more so that
\begin{equation}\label{eq:reader-common-error-budget}
 x^{-1/8}+p_x\le\varepsilon_f(t),\qquad
 x^{-1/32}\le\varepsilon_f(t),\qquad
 x^{-1}\le\varepsilon_f(t)^2
 \quad(t/2\le x\le t).
\end{equation}
For the first inequality apply \eqref{eq:reader-band-power} with
$c=f+1$ and note that
$2\varepsilon_{f+1}(t)\le\varepsilon_f(t)$ eventually. For the
second use $c=f,v=1/32$; for the third use $c=2f,v=1$ and
$\varepsilon_{2f}=\varepsilon_f^2$.

Thus one $U$ is chosen first, then one entrance threshold, then
for each $x$ one $p_x$, and all the estimates hold for every
admissible coupling pair, period, forest and tag. No maximum over
an infinite collection of periods is being taken: each input
estimate already has its uniform threshold.

\subsection{The tube and covariance belong to the retained history}

Let $C$ be the selected family and fix its actual precision horizon
and path attachment. At an adjacent step $i<N$, form the local
Gaussian data from the fine path at $i+1$, the coarse path at $i$,
their prescribed clips, and the tower between those very levels.
The path bounds come from the Wilson representations of $C$;
the metric floors, path attachments, symmetry and crossing
identities come from this horizon. They are therefore the same
objects in the local Gaussian estimate and in the short-covariance
construction. Replacing either by another witness with similar
bounds would not justify the calculation below.

Suppose the step is in a constant-order plateau,
$J(i+1)=J(i)$. Its strip width, metric floor and price then have
the same stage index on both sides. Write $P_i^{\rm cond}$ for
its conditional curvature covariance, and $Q_i^{\rm short}$ for
the prescribed finite-range covariance used in the reader
expansion. The common-cover estimate, applied to the actual fine
metric, gives
\begin{equation}\label{eq:reader-history-short-error}
 |P_i^{\rm cond}(a,b)-Q_i^{\rm short}(a,b)|\le p_x.
\end{equation}
This includes the two cases in the short-row rule. Where the rule
keeps the covariance entry there is no discrepancy. Where it
sets an entry to zero, the identically truncated cover row has
zero at that periodic address, so its tail estimate bounds the
original entry. All images in the periodization are included in
that support test.

For a posed actual sector, the local estimate gives
$|C_i^{\rm read}(a,b)-P_i^{\rm cond}(a,b)|\le x^{-1/8}$ on
both reader supports. Adding \eqref{eq:reader-history-short-error}
and using \eqref{eq:reader-common-error-budget} yields
\begin{equation}\label{eq:reader-selected-sight}
 |C_i^{\rm read}(a,b)-Q_i^{\rm short}(a,b)|
 \le\varepsilon_f(t).
\end{equation}
The short covariance is a comparator. The random vector continues
to have the actual sector covariance $C_i^{\rm read}$.

The corresponding transport-row estimate uses the same horizon
and gives the short mean-transport error. Combining these two
estimates on each plateau step is the history comparison used by
the reader expansion. The common schedule must place the actual
entrance parameter in $[t/2,t]$ and the actual couplings in
\eqref{eq:reader-coupling-band}; these are the purchased plateau
and coupling conditions retained in the application.

\subsection{One pose contains all required reads}

We give the geometric check, since it prevents accidental use of
different local gauges for different readers. Put
$b_x=\lceil(\log x)^2\rceil$ and $w_x=b_x^2$. At a fixed
step let $y$ be the next prescribed scale parameter, and suppose
the paid room conditions hold. In explicit terms the conditions
used here are
\[
 b_x\ge200,\quad w\ge d+10,\quad
 w w_y+L\le Lw_x,
\]
together with room nonwrapping at both levels and the prescribed
reader separation of at most $d w_x$. Here $d$ is the fixed
reader separation order. The recording parameter $r(x)$ satisfies
\[
 x\le r(x)\le x^2,\qquad b_{r(x)}\le4b_x,
\]
and its flag threshold equals the healthy threshold at $x$.
Consequently $11b_{r(x)}+1\le w_x$. Also the nonwrapping room
bound at radius $ww_x$ implies the local Gaussian side bound:
$100(1+b_{r(x)})\le2(ww_x+1)$ eventually under these
inequalities, and the torus side exceeds the latter quantity.

For this application fix the seed-free store: both its seed sets
are empty, and its flag set is exactly the set of flags of the
actual coarse field at recording threshold $\operatorname{temp}(r(x))$.
This threshold equals $\tau_x$. Thus, for a coarse field healthy
on the whole room of radius $ww_x$, the store has no marked site
in that room. Health alone would not remove seeds from an
arbitrary store. For this seed-free store,
the local axial construction gives one gauge on this room and a
logarithmic direction of norm at most $8ww_x\tau_x$, where
$\tau_x$ is the healthy threshold. The purchased chart conditions
place this inside half the logarithmic clearance allowance and
inside the fixed chart radius. Extend that direction by zero
outside the room. It remains one global direction with the same
supremum bound. At any address whose distance from the room
boundary is at least $w_x$, the whole $b_{r(x)}$-ball lies
inside the chart, and its $11b_{r(x)}+1$ seed neighbourhood
contains no flag. The replica bound follows from the linear
replica norm bound and the exponential chart estimate. These are
exactly the pose conditions used above.

It remains to fit the fine guards into that room. A fine bond in
the enlarged guard room of radius $ww_y+8$ has parent at distance
at most $w_x+8$ from the parent of the fine reader centre.
Indeed the parent-distance inequality is
\[
 L\,\operatorname{dist}(\operatorname{parent}u,
                         \operatorname{parent}v)
 \le\operatorname{dist}(u,v)+L,
\]
and $ww_y+L\le Lw_x$ gives the assertion. The coarse centre
of a tested reader is at distance at most $dw_x$ from the
chosen root reader. Adding distances shows that every such guard
bond has
\[
 \operatorname{dist}(o,\operatorname{parent}v)+w_x+2
 \le(d+2)w_x+10\le ww_x,
\]
using $w_x\ge8$ and $w\ge d+10$. Every incidence read has
the same inequality with $+1$ in place of $+2$. Thus both
endpoints of every covariance entry, and all four bonds used in
every nonlinear read, lie in the one pose. Signed translations
carry the reader centres, masks, forest and tag together, so the
same inequalities apply to every selected orientation.

\subsection{The joint reader model on the actual sample space}

Fix a forest and tag, retain this same seed-free store, and choose
a chart cutoff $\delta>0$ with $2\delta$ smaller than the fixed
first logarithmic chart radius of the gauge group. These are the
cutoff hypotheses of the measurable field construction in the
joint reader model. Use the common product space of frozen Haar
links and the centered Gaussian sector coordinates from
Section~\ref{sec:ex-lattice-all-arity-reader-gaussian-model}.
The signed transport and the one pose gauge send the store,
forest, tag and sample together. Their measure-preserving
identities give one joint centered incidence vector, with
covariance $C_i^{\rm read}$, for every reader. Equation
\eqref{eq:reader-selected-sight} is therefore an entrywise
comparison of that actual joint law with the short covariance,
including all off-diagonal pairs.

For each reader use the guard on its enlarged fine room. The
Gaussian guard estimate gives
\[
 0\le\Gamma_a\le1,\qquad
 \int(1-\Gamma_a)\,d\mu\le x^{-1}
                          \le\varepsilon_f(t)^2.
\]
On its positive set, the nonlinear comparison makes the fine
field healthy on that reader patch and gives one gauge element
for the whole patch. In that gauge, each normalized coordinate
differs by at most $x^{-1/32}\le\varepsilon_f(t)$ from the
sum of the fixed parent forecast and the corresponding coordinate
of the common Gaussian vector. The forecast uses the actual
coarse coupling $\beta'$; the fine field uses $\beta$. Their
conversion was included in the adjacent coupling estimate, so
they are not set equal in this step.

The same construction works for every permitted reader in the
fixed finite family. Its mask bound is per reader, with no union
over the family, and its gauge is common to the coordinates of
one patch. All readers still share the same random vector and
the same parent forecast. These are precisely the quantitative
receipts left explicit in the joint reader-model construction.
They are obtained on the selected family whenever its purchased
plateau, room, chart and coupling conditions hold. The later
activity reassembly remains responsible for identifying the
resulting expansion with the interacting block integral.
